\section{Introduction, Literature Review, and Summary}

In canonical models of firm investment, firms are perfectly rational agents capable of 
solving infinite-horizon dynamic optimization problems. The firm uses capital $k_{j,t}$ 
(physical machines, equipment, and infrastructure) to produce output via a production 
function $f(z_{j,t}, k_{j,t})$, where $z_{j,t}$ is a stochastic productivity shock that 
captures fluctuations in the firm's efficiency (e.g., demand shifts, supply 
disruptions, technological changes). Each period, the firm can purchase new capital by 
investing $i_{j,t}$, but the firm incurs convex adjustment costs $\psi(i_{j,t}, k_{j,t})$ that penalize rapid changes to its capital 
stock (frictions of retooling factories, training workers on new equipment, and reorganizing production). 
Capital depreciates over time at rate $\delta$ (machines wear out, equipment becomes obsolete), 
so the firm's capital stock evolves as $k_{j,t+1} = (1-\delta)k_{j,t} + i_{j,t}$: what 
survives from today plus what is installed. To finance investment, the firm can issue bonds $b_{j,t+1}$ (borrow from 
external creditors), but lenders require collateral: borrowing is capped at a fraction 
$\theta$ of next-period capital.

This problem admits an analytical closed-form solution.
Under this so-called neoclassical model, the optimal investment rule is simple: invest 
until the shadow value of an additional unit of capital (Tobin's $q$) equals its 
installation cost. When $q > 1$, a new machine generates more value than it costs to 
install, so the firm invests; when $q < 1$, existing machines are worth less than their 
replacement cost, so the firm divests.

While elegant, there are well-documented empirical irregularities that pose challenges
to this model.
\begin{enumerate}
    \item \textbf{Weak investment-q relationship.} The neoclassical model predicts that 
    Tobin's $q$ is a sufficient statistic for investment. Empirically, $q$ explains almost 
    none of the variation in firm-level investment. \citep{FazzariHubbardPetersen1988,Hubbard1998,EricksonWhited2000, AbelEberly1996}.
    \item \textbf{Excess sensitivity to cash flow.} The neoclassical model predicts 
    investment depends only on $q$, not on internal funds. Empirically, investment responds 
    strongly to cash flow even after controlling for $q$. \citep{FazzariHubbardPetersen1988,KaplanZingales1997,Lamont1997, AbelEberly1996}.
    \item \textbf{Excess investment dispersion.} The neoclassical model predicts firms 
    with identical fundamentals make identical choices. Empirically, firms with similar 
    observables make persistently different investment decisions. \citep{HsiehKlenow2009, BachmannBayer2014}
    \item \textbf{Investment inertia and lumpy adjustment.} The neoclassical model with 
    convex costs predicts smooth, continuous adjustment. Empirically, firms underreact for 
    extended periods then make large discrete changes. \citep{DomsDunne1998,CooperHaltiwanger2006}.
    \item \textbf{Asymmetric investment dynamics.} The neoclassical model predicts 
    symmetric responses to positive and negative shocks. Empirically, firms expand gradually 
    but contract sharply.\citep{AbelEberly1996,CaballeroEngel1999}.
    \item \textbf{Fat-tailed firm size distribution.} The neoclassical model predicts 
    convergence to a common steady state. Empirically, firm sizes exhibit persistent 
    heterogeneity with a fat right tail. \citep{Lucas1978, Luttmer2007}
\end{enumerate}

This project proposes to address these puzzles by adapting the framework of 
\cite{IlutValchev2023}\footnote{\cite{IlutValchev2023} develop their framework in the 
context of a consumption-savings problem. This project extends it to the richer and 
empirically distinct setting of firm investment.} to the context of firm investment. In 
this model, agents do not know the true value of different actions. They may hold some 
initial beliefs (a prior) over the action-value function $Q^*(s, a)$, which encodes the 
present discounted value of taking action $a$ in state $s$ and behaving optimally 
thereafter, but they must learn and refine these beliefs over time through experience and 
reasoning (planning). From the firm's perspective, the true action-value function $Q^*$ is 
unobserved, and the firm's beliefs $(\hat{Q}_t, \hat{\Sigma}_t)$ about it become part of 
the relevant state.

The firm operates as an RL agent that explores the state-action space to learn 
about $Q^*$, it uses temporal-difference learning (\textit{model-free} RL) to update 
its beliefs as well as some planning (or reasoning) (\textit{model-based} RL) to refine 
them, and it ultimately develops an investment policy to maximize the expected present value of dividends.
The true underlying MDP is, in fact, the neoclassical model described above, so the true 
$Q^*$ is the known solution to that problem. However, firms in exploring and learning
about $Q^*$ may generate a dispersion of investment policies that are path dependent as
they are each subject to idiosyncratic shocks. This dispersion of policies may be persistent and
not perfectly aligned with the neoclassical solution. In this framework we generate an 
ergodic distribution of firm-level investment policies that can match the empirical puzzles.



% \paragraph{Related Literature.}
% This paper contributes to several strands of the literature. Most directly, it builds
% on \cite{ilutvalchev2023}, who develop the GP learning and rational inattention
% framework for household consumption and show that it generates learning traps,
% persistent hand-to-mouth behavior, and high marginal propensities to consume. We
% extend their framework to the firm side, where the forward-looking nature of capital
% accumulation, the durability of the state variable, and the relevance of firm
% heterogeneity for aggregate productivity create distinct economic forces.

% On the investment side, the paper relates to the literature on lumpy investment and
% non-convex adjustment costs \cite{cooperhaltiwanger2006, khanthomas2008}. Our model
% generates observationally similar patterns from a different
% primitive---informational rather than technological. \cite{bloom2009} and
% \cite{bloometal2007} emphasize uncertainty shocks as a driver of investment
% fluctuations through real options effects. In our framework, uncertainty is not an
% exogenous shock but an endogenous state that evolves with the firm's learning history.

% The paper contributes to the literature on misallocation and productivity.
% \cite{hsiehklenow2009} and \cite{restucciarogerson2008} document and quantify the
% aggregate productivity losses from cross-sectional dispersion in marginal products.
% Subsequent work has attributed this dispersion to adjustment costs
% \cite{askeretal2014}, financial frictions \cite{midriganxu2014, moll2014}, and
% information frictions \cite{davidhopenhaynvenkateswaran2016}. Our learning trap
% mechanism is closest to the information frictions channel but differs in that the
% persistence of misallocation arises endogenously from the optimal choice to stop
% learning, rather than from exogenous signal noise.

% More broadly, the paper connects to work on bounded rationality and firm
% decision-making. \cite{sims2003} introduces rational inattention, and
% \cite{gabaix2014} develops a sparsity-based model in which agents attend to only a
% subset of relevant variables. \cite{malmendiernageletal2011} and
% \cite{malmendiertateyan2011} provide empirical evidence that managers' personal
% experiences persistently shape corporate financial decisions, consistent with the
% experience-based learning mechanism in our model. \cite{fusterlabsonmendel2010}
% study how extrapolation from recent experience generates excess macroeconomic
% volatility. Our contribution is to embed these ideas in a formally disciplined
% framework with provable ergodic properties, connecting the micro evidence on
% experience effects to the macro evidence on investment dynamics and misallocation.